% Generated draft: requires mathematical and source review.
\documentclass[11pt]{article}
\usepackage[margin=25mm]{geometry}
\usepackage{fontspec}
\setmainfont{Latin Modern Roman}
\usepackage{amsmath,amssymb,mathtools,unicode-math}
\setmathfont{Latin Modern Math}
\usepackage{enumitem,needspace,xurl,hyperref}
\setlength{\parindent}{0pt}
\setlength{\parskip}{5pt}
\setlength{\emergencystretch}{3em}
\newcommand{\N}{\mathbb{N}}
\newcommand{\Z}{\mathbb{Z}}
\newcommand{\Q}{\mathbb{Q}}
\newcommand{\R}{\mathbb{R}}
\newcommand{\C}{\mathbb{C}}
\newcommand{\F}{\mathbb{F}}
\newcommand{\A}{\mathbb{A}}
\newcommand{\PP}{\mathbb{P}}
\newcommand{\E}{\mathbb{E}}
\newcommand{\eps}{\varepsilon}
\newcommand{\dd}{\,\mathrm{d}}
\begin{document}
\section*{UnsolvedMath ID: 30002262}
\textbf{Source problem code:} catalogue entry
\subsection{Definitions and mathematical statement}
Let \(\mathcal L\) be a first-order language extending the language of rings, and let \(K\) be an infinite field equipped with an \(\mathcal L\)-structure extending its field structure. Call \(K\) \emph{minimal in \(\mathcal L\)} if, for every \(\mathcal L\)-formula \(\varphi(x,\bar y)\) with one field variable \(x\) and every finite tuple \(\bar a\) of parameters from \(K\), the set
\[
\{x\in K:K\models\varphi(x,\bar a)\}
\]
is finite or has finite complement in \(K\). This uses the usual convention that a minimal field is infinite and allows parameters; it does not require every elementary extension to be minimal. The supplied record does not specify the language, so \(\mathcal L\) remains an explicit formulation parameter.

\textbf{Algebraic-closure question (attributed to Podewski's conjecture).} Is every minimal field algebraically closed? For a fixed intended language \(\mathcal L\), the proposed assertion is
\[
\forall K\;\left(K\text{ is a minimal field in }\mathcal L\ \Longrightarrow\ \forall f\in K[X]\setminus K\;\exists a\in K\;(f(a)=0)\right).
\]
The question ranges over all characteristics; the characteristic-zero restriction in the record belongs to the order question below.

\textbf{Order question.} Can a minimal ordered field of characteristic zero exist? Equivalently, the original wording asks whether there are no such fields. Under the usual algebraic meaning of an ordered field, this asks whether there exists a pair \((K,<)\) such that \(K\) is minimal in the intended language \(\mathcal L\), \(\operatorname{char}(K)=0\), and \(<\) is a strict total order satisfying
\[
\begin{aligned}
&\forall a,b,c\in K\quad (a<b\Longrightarrow a+c<b+c),\\
&\forall a,b\in K\quad (0<a\ \land\ 0<b\Longrightarrow 0<ab).
\end{aligned}
\]
This is a conditional interpretation of the wording: the record does not specify whether the order is external to the structure, definable in it, or included in \(\mathcal L\), nor whether the cited discussion uses a different model-theoretic order convention. These possibilities must not be silently identified.

\textbf{Unresolved original component.} The original extraction also asks whether an unspecified constructed group is quasi-minimal, or at least regular. No construction or identifying data for that group, and no intended definition of regularity, are supplied. Consequently that component cannot be given a precise mathematical statement here. The supplied statement-verification note says that the report separately poses the two field questions and that the group component was omitted from the corrected statement because its defining context was insufficient; this is supplied provenance, not an independent verification.
\subsection{Short English statement}
The field questions ask whether every infinite minimal field is algebraically closed and whether a minimal ordered field of characteristic zero can exist. Minimality means that every definable subset of the field, allowing parameters, is finite or cofinite. The intended language and meaning of the order need clarification. The original group question also remains unformulated because its construction is missing. This is a statement-only draft requiring human mathematical review.
\subsection{Sources}
\begin{itemize}
\item Supplied local problem-registry record for permanent UnsolvedMath ID \(30002262\), source code \texttt{catalogue entry}, entitled \emph{Algebraic Closure and Orderability of Minimal Fields}, in the Oberwolfach Reports Open Problems collection. The project registry is \texttt{data/unsolvedmath/problems.json}, with additional records in \texttt{catalog/problems.json}; no independent inspection of either file is asserted. The supplied proposed year is 2013; no proposer is named. The algebraic-closure conjecture is attributed to Podewski.
\item \emph{Model Theory: Groups, Geometry, and Combinatorics} (2013), Oberwolfach Reports, Report 01/2013, pp.\ 5--66, DOI \texttt{10.4171/owr/2013/01}: \url{https://doi.org/10.4171/owr/2013/01}. This is the supplied source report and is also listed among the supplied literature sources. The background describes these reports as records of mathematical workshops and their open-problem discussions.
\item \emph{Around Podewski's Conjecture}, supplied literature reference, DOI \texttt{10.4064/fm222-2-4}: \url{https://doi.org/10.4064/fm222-2-4}. No author or further bibliographic details are supplied.
\item Literature assessment dated exactly \texttt{2026-08-21}: the supplied status is \texttt{partially\_solved}. The supplied assessment states that generic stability proves algebraic closedness for regular fields and resolves broad special cases, while the constructed group and remaining minimal and quasi-minimal questions are not fully settled. This is supplied dated curation derived from the OWR material, not an independently verified literature assessment. Its time-sensitive status and cited mathematical claims require independent verification.
\item The supplied statement status is \texttt{corrected\_verified}. Its verification note asserts that the official report separately poses the two field questions and that the extraction's additional constructed-group question lacks enough defining context to publish safely. This label and explanation are preserved as source metadata; no source visit or independent mathematical verification is claimed.
\end{itemize}
\end{document}
